\section{Discussion}
\textbf{What the agents add.} Beyond the earlier hub, the contribution that can be defended is structural: two stateful agents take real decisions (difficulty proposals, probes, coaching plans) through typed messages, with an authority matrix enforced in code and arbitration that is transparent when the agents contribute nothing (identical trajectories with no usable agent). The Strategy Agent's ledger is what a threshold rule or PPO's 6-D state lacks; the replay shows that it changes the difficulty path in every session, but \emph{not} whether the change is better.

\textbf{What the safeguards do.} The ablations separate roles. Schema validation contains malformed and forbidden output; the authority check adds classification; the confidence threshold prevents ungated ``steady'' proposals from overriding PPO (isolated in XC; the removal of arbitration in XA changes four checks at once); guardrail precedence prevents bypass; grounding prevents ungrounded emphasis; the per-agent breaker bounds repeated failures at the cost of availability. None of these detects a well-formed, confident, wrong proposal, which was followed in a measurable share of decisions; only bounds and guardrails limit its effect.

\textbf{Why two agents rather than one rules engine?} The evidence does not show that two agents are needed. Removing the Coaching Agent changes no difficulty, score or question path (A5), so one merged component would have behaved identically on these metrics. The split is an authority and state boundary: each agent has its own private ledger, its own typed output and its own breaker, so a fault in one does not remove the other (XB), and coaching can be audited and disabled without touching difficulty. Whether this modularity is worth its overhead for a system this size, against a single rules engine with the same ledger, was not tested.

\textbf{Is it multi-agent?} By the weak notion of agency~\cite{wj95}, each agent has private persistent state, reacts to observations, makes pro-active proposals (a probe the thresholds would not trigger) and communicates (Strategy to Coaching through the coordinator). Their decision logic is rule-based; they are not LLM agents, do not negotiate, and the system remains centrally orchestrated. This is bounded-authority role specialisation, not decentralised intelligence, and PPO's role is unchanged.

\textbf{Relation to the tracks.} The work fits agent-native development and lifecycle engineering (typed contracts, authority, fault containment, regression-tested integration) better than core multi-agent collaboration; the authority matrix, containment of faulty agents and breaker behaviour also relate to the security-and-resilience track. We have not verified the official call text (Sect.~\ref{sec:limits}).
